\chapter{Черга після інтеграції translation system/local TeX worker-модулів}
\label{chap:session07-queue}

Після інтеграції повного translation system source package порядок роботи змінився. Тепер не треба витрачати глибокі токени на те, що вже зроблено як source-preserved модуль. Потрібно витрачати їх на контроль якості, формули, терміни, і вузли, які склеюють корпус.

\section{P0: найближче}

\begin{enumerate}[leftmargin=*]
\item \textbf{Solà ESKF full paper.} Завершити source-preserved модуль: quaternion conventions, error-state, IMU kinematics, noise, injection/reset.
\item \textbf{Micro Lie full paper.} Використати local TeX worker fragments як старт, але пройти весь paper file-by-file з компіляцією.
\item \textbf{PySDR Ukrainian RST-to-TeX.} Це дешевий виграш: український RST уже існує, потрібне машинозчитуване TeX-представлення.
\item \textbf{Labbe Kalman.} Конвертувати перші практичні notebook/markdown глави у TeX: g-h filter, discrete Bayes, Gaussian, Kalman predict/update.
\item \textbf{Correll repair.} Залишити шість розділів як standalone book module; виправити refs, figure placeholders, captions, source labels.
\end{enumerate}

\section{P1: broad applied-math balance}

\begin{enumerate}[leftmargin=*]
\item FDM/FEM: wave/heat/Poisson, CFL, convergence, manufactured solutions.
\item Open optimization: least squares, convexity, LP/QP, robust losses, constrained optimization.
\item Control systems: state-space, controllability/observability, PID/LQR, discrete-time models.
\item Circuits/RF basics: filters, impedance, transmission lines, noise, link budget as education-only material.
\end{enumerate}

\section{Public/local split}

Public: open-source or source-backed educational TeX, compiled PDFs, terminology, build scripts, QA notes.

Local review/staged: scan OCR drafts, source-repair fragments, incomplete formula reconstructions, and materials that need provenance or mathematical QA before being treated as reader-ready.
